\section{Preliminary Analysis as a Diagnostic Tool for Multimodal Assessment}
The preliminary analysis aims to make the MAIA assessment more diagnostic and less dependent on final task outcomes. Performance on a Visual Question Answering (\acs{VSV}) item alone does not reveal which prerequisite capabilities required by the task have been successfully satisfied. A correct answer does not necessarily imply that the model has accurately identified the relevant entities, reconstructed the events, established the required temporal relations, or correctly integrated the available evidence. Conversely, an incorrect answer does not indicate which of these components is responsible for the failure.
This limitation is particularly relevant for MAIA, which evaluates fine-grained reasoning grounded in video content. For this reason, the preliminary analysis is performed independently of the \acs{VSV} task. Rather than relying exclusively on direct question answering and aggregated accuracy, the video content is first represented through a structured analysis of entities, events, and relations among them. This representation is subsequently compared with an analogous analysis of the requirements imposed by the questions. Separating the two stages reduces the risk that video analysis is conditioned by the linguistic formulation of the question or by the plausibility of the available answer options. The objective is therefore to determine \textbf{what information can be retrieved from the video and what level of integration is required to retrieve it}, before relating this evidence to the demands of the \acs{VSV} task.
This framework has a methodological affinity with \textit{VILMA} \cite{kesen2024vilma}, particularly with its distinction between preliminary proficiency evaluation and the main test. \textit{VILMA} implements this principle through a caption--foil paradigm in a zero-shot setting, using proficiency tests to verify prerequisite capabilities before interpreting performance on more complex phenomena. Our approach does not reproduce the \textit{VILMA} evaluation format; rather, it adopts the underlying methodological principle that complex abilities should be interpreted in relation to the simpler perceptual, temporal, and semantic capabilities on which they depend.
Accordingly, the proposed analysis distinguishes between tasks that require direct retrieval of localized information and tasks that require progressively greater integration across entities, events, temporal contexts, and inferential relations. In video understanding, this distinction is particularly important because relevant information may be contained within a single frame or distributed across a sequence. Questions belonging to the same nominal category may therefore impose substantially different processing requirements. A spatial question, for example, may require the retrieval of a location from a single frame or the reconstruction of a spatial configuration after a sequence of events. To capture this variation systematically, the classification is based on explicit operational criteria rather than on an intuitive notion of difficulty.
The proposed question classification framework estimates the complexity of the operations required to retrieve and integrate the relevant evidence. A common three-level scale is applied to three dimensions: \textbf{spatial}, \textbf{temporal}, and \textbf{causal}. At Level~0, the required information can be retrieved directly without comparison or reconstruction. Level~1 requires the integration of an explicit relation or dependency involving a limited amount of evidence. Level~2 requires a broader reconstruction across events, temporal phases, state changes, or implicit relations. The same progression is applied consistently across the three dimensions, as summarized in Table~\ref{tab:question_classifications}.
The classification criteria are grounded in properties of the evidence required to answer each question, including its temporal distribution, the number of relevant evidence units, the presence of event relations, and the explicit or implicit nature of the required information. These properties are intended to describe the informational requirements of the task rather than its linguistic formulation. An independent analysis of the video is therefore necessary to determine whether the entities, events, and relations presupposed by the question are actually recoverable from the available visual evidence.
The video analysis is organized into successive stages of semantic representation. The preprocessing stage provides temporally indexed video content, including a dense representation used to preserve temporal information. Four omnimodal and multimodal models, subsequently evaluated on the \acs{VSV} task, independently analyze temporally localized portions of each video. Their outputs are used to construct structured representations of the entities, events, and relations identified in the observed content.
The analysis is performed independently of the questions so that the resulting representation reflects information retrieved from the video rather than evidence selected in response to a specific query. Temporally overlapping portions of the video are used to preserve the connection between extracted information and the intervals in which it is supported. This temporal anchoring is necessary for the subsequent comparison between the evidence represented by the models and the evidence required by individual questions.
For each video, the analysis progressively considers three levels of representation: \textbf{entities}, \textbf{events}, and \textbf{relations between events}, including inferential relations where required. Outputs from the different models are initially kept separate, allowing inter-model agreement to be measured without prematurely merging their predictions. Agreement does not constitute ground truth, since different models may produce the same unsupported interpretation; it is instead treated as an indicator of \textbf{prediction stability}. The analysis therefore distinguishes information that is directly supported by observable evidence from information introduced through model inference.
Within this framework, \textbf{event extraction} is treated as the identification of semantically structured events associated with specific temporal intervals. Because the same event may appear in multiple overlapping portions of the input, subsequent consolidation reduces redundancy while preserving genuinely repeated events. The consolidated representation can then be used to establish temporal relations such as \textit{before}, \textit{after}, \textit{during}, and \textit{overlap}, producing a temporally organized representation of the video.
This distinction between entities and events is essential for evaluating temporal complexity. Knowing that a video contains a person, a mug, and a table provides a different type of information from knowing that the person picks up the mug, drinks from it, places it on the table, and subsequently walks away. The former represents the entities present in the scene, whereas the latter requires the representation of actions, temporal order, and changes of state. Similarly, identifying the location of an object at a given moment imposes different requirements from identifying its location after a particular event. The preliminary analysis is designed to make these differences explicit.
The causal dimension introduces an additional inferential requirement. Temporal succession alone does not establish causality; consequently, causal relations must be distinguished from simple co-occurrence or temporal ordering. The analysis first represents the relevant events and their temporal organization. It can then be determined whether answering a causal question requires information that is explicitly available in the video or an additional inference connecting observed events, states, or intentions.
A further objective of the preliminary analysis is to determine which information can be recovered from the visual modality independently of other channels. This addresses a methodological issue highlighted by previous experiments \cite{deodati2026lost}, in which visual information was compared with additional modalities such as audio and automatic transcription. In the original MAIA setting, the addition of these channels did not produce a systematic improvement, while the inclusion of audio transcriptions resulted in a small but statistically robust performance decrease.
One relevant characteristic of the original dataset is that its \acs{VSV} items were constructed around visual content rather than around information specifically provided by the acoustic channel. For this reason, the preliminary analysis establishes a visual baseline focused on \textit{what can be seen} before assessing the contribution of \textit{what can be heard}. This separation makes it possible to evaluate whether acoustic information provides genuinely complementary evidence and whether questions designed for multimodal assessment require information distributed across more than one modality.
The purpose of the preliminary analysis is therefore to characterize \textbf{which visual information is available in the video and how much evidence must be retrieved and integrated to answer a question}. This representation provides the basis for evaluating whether the models subsequently tested on the \acs{VSV} task are able to recover and combine the information required by each item.
In this regard, the preliminary analysis moves \textit{MAIA-AV} beyond a purely result-based evaluation towards a more diagnostic and competence-oriented framework. Building on the methodological principle underlying the proficiency tests introduced by \textit{VILMA}, the proposed approach makes the prerequisite conditions for successful question answering explicit through a structured analysis of video content. This allows model performance to be interpreted in relation to specific task requirements, distinguishing limitations in entity recognition, event identification, temporal integration, spatial reasoning, and causal inference rather than attributing an incorrect response generically to a failure of multimodal reasoning.
\begin{table}[htbp]
    \centering
    \small
    \renewcommand{\arraystretch}{1.4}

    \begin{tabularx}{\textwidth}{|l|X|X|X|}
        \hline
        \rowcolor{EcoForestGreen}
        \textcolor{white}{\textbf{Category}} &
        \textcolor{white}{\textbf{Level 0}} &
        \textcolor{white}{\textbf{Level 1}} &
        \textcolor{white}{\textbf{Level 2}} \\
        \hline

        \cellcolor{EcoForestGreen}
        \textcolor{white}{\textbf{Spatial (S)}} &
        Direct retrieval of an entity's location or spatial property from a single frame. &
        Explicit spatial relation or event-conditioned localization involving at most two evidence units. &
        Spatial reconstruction across multiple phases, including movement, relocation, or state change. \\
        \hline

        \cellcolor{EcoForestGreen}
        \textcolor{white}{\textbf{Temporal (T)}} &
        Recognition of a single event or state without temporal comparison. &
        Temporal localization or pairwise ordering between two explicit events. &
        Multi-event temporal reconstruction, including duration, repetition, onset, or state evolution. \\
        \hline

        \cellcolor{EcoForestGreen}
        \textcolor{white}{\textbf{Causal (C)}} &
        Retrieval of an explicitly stated cause without inference. &
        Explicit cause--effect relation involving at most two evidence units. &
        Implicit or multi-event causal reconstruction requiring contextual or intentional inference. \\
        \hline
    \end{tabularx}

    \caption{Complexity levels for spatial, temporal, and causal questions.}
    \label{tab:question_classifications}
\end{table}